\chapter{Working with survey platforms}\label{ch:surveys}

The survey closes Friday. It is Tuesday, the program director asks how response is going, and nobody can answer without logging into a web dashboard and eyeballing a number that no one wrote down. Survey platforms hold your data behind their websites, and the evaluator who can only reach it by hand is an evaluator who cannot monitor, cannot reproduce, and cannot move fast.

This chapter connects Stata directly to the survey platforms evaluators actually use, so responses download themselves, response rates can be watched while the survey is live, and the instrument's meaning travels with its data across waves. Automation here is not a convenience: until the pull is scripted, there is nothing to monitor. Chapter~\ref{ch:workshop} ruled out spreadsheet analysis for the same reason; this is that rule, applied to the survey.

Behind everything in this chapter sits one framework, total survey error\index{total survey error}: the accounting of every gap between the estimate a survey reports and the quantity it means to measure. \textcite{groves2010tse} trace how the concept grew from a tool for reasoning about sampling into a full ledger of coverage, nonresponse, and measurement error, and \textcite{biemer2010tse} turns that ledger into a design discipline that spends the budget where the largest error lives. For an embedded evaluator the payoff is concrete rather than theoretical: when the monitor of Section~\ref{sec:monitor} shows a rate falling, the ledger names what the fall threatens, and the nonresponse diagnostics later in the chapter test whether the threat came true, since a falling rate is not a problem in itself but a warning that one term of the error ledger may be growing. The stakes are not academic. \textcite{meyer2015surveys} document three decades of rising nonresponse, imputation, and misreporting in the household surveys that anchor U.S. policy, and an evaluator meets that same erosion in miniature every time a program survey's response rate slips. Watching that rate while the survey is open, and knowing which term of the ledger it threatens, is the unglamorous core of longitudinal practice.

Two design commitments frame everything that follows. The first is tailored design: \textcite{dillman2014tailored} show that response is built at the instrument, through mode, contact sequence, and question wording, long before any download, so the automation in this chapter is stewardship of an instrument someone designed with care, not a substitute for that care. The second is a shared vocabulary for what a response rate even is. \textcite{aapor2023standard} fix the case dispositions and outcome-rate formulas that let one evaluator's ``62 percent'' mean the same thing as another's; without them, a funder comparing this wave's rate to last year's is comparing a number that quietly redefined itself along the way.

\section{Pulling responses automatically}
Make your responses download themselves, so the data that lives behind a vendor's login becomes a file you can rebuild on command. Every major platform exposes an API\index{API!survey platforms}, and the patterns differ mainly in how much ceremony each demands. REDCap\index{REDCap} is the simplest: a single POST with \texttt{content=record\&format=csv} returns the records, which is why server-hosted REDCap, friendly to the institutional review board that must approve how human-subjects data is handled, is straightforward to script. Qualtrics\index{Qualtrics} uses a three-step export, request the file, poll until it is ready, then download, because large exports are prepared asynchronously.\marginnote{The auth models differ too: REDCap issues one project-scoped API token that both identifies you and names the dataset, while Qualtrics wants an account-wide \texttt{X-API-TOKEN} header plus a separate datacenter ID and survey ID. One REDCap secret, three Qualtrics coordinates.} SurveyMonkey pages through responses in bulk but rate-limits hard and reserves full export for paid tiers, a friction worth knowing before you promise a client a live feed. KoBoToolbox, common in field and humanitarian work, returns responses as JSON from an asset endpoint.

A first pull deserves a small rehearsal. Before scripting the full download, pick two respondents you can identify on the platform's dashboard, pull only their records, and check the export against the screen; multi-selects, timestamps, and anything with an accented character are where the export and the dashboard most often disagree. Ten minutes on two known cases settles the export options while the stakes are still zero.

The quickest payoff needs no platform API at all. If a Google Form's response sheet is link-shared, one line pulls it into Stata:
{\footnotesize
\begin{verbatim}
local sheet "https://docs.google.com/spreadsheets/d/KEY"
import delimited "`sheet'/export?format=csv&gid=0", clear
\end{verbatim}
}
Hand-downloading is the survey version of spreadsheet analysis: it works once and reproduces never. Scripting the pull makes the response data replicable and, as the next section shows, makes live monitoring possible in the first place.

\subsection{Recombining multi-select exports back into one variable}
The pull is only the first surprise the platform hands you. Qualtrics and Google Forms export a ``check all that apply'' question not as one variable but as a block of one-hot indicator columns, one per option, so a five-option question arrives as five \texttt{0/1} variables named after the choices. That layout is right for a check-all item, where a respondent can pick several, but wrong for a single-answer question that the platform still splits, and wrong for the analyst who wants one labeled categorical to tabulate. Rebuilding it by hand, an \texttt{egen group()} followed by a stack of \texttt{cond()} and \texttt{label define} lines, is the small error-prone chore that \texttt{undummy}\index{undummy@\texttt{undummy}} does in one line. It is the exact inverse of \texttt{tabulate, generate()}, which is how survey platforms produced the dummy block in the first place.

\begin{kaobox}[frametitle=Package Integration: \texttt{undummy}]
\textbf{Integrated command:} \texttt{undummy} recombines a set of mutually exclusive dummy variables into one labeled categorical, mapping labels from the value each dummy carries (or, with \texttt{varnames}, from the dummy names themselves). Install with the house idiom:
{\footnotesize
\begin{verbatim}
local gh "https://raw.githubusercontent.com/ericabooth"
net install undummy, ///
    from("`gh'/undummy-stata-public/main/") replace
\end{verbatim}
}
The inverse of \texttt{tabulate, generate()}. On the classic \texttt{auto} data the round trip is one split and one recombine:
{\footnotesize
\begin{verbatim}
sysuse auto, clear
tab foreign, gen(fd)
undummy fd1 fd2, generate(origin) varnames
tabulate origin
\end{verbatim}
}
The new \texttt{origin} carries \texttt{r(k)=2} categories labeled from the dummy names, and \texttt{return list} reports \texttt{r(base)}, the first dummy, and \texttt{r(generate)}, the new variable's name. Two caveats govern honest use. First, run \texttt{undummy, checkdummies} on a survey block before you trust it: a genuine check-all-that-apply set is not mutually exclusive, and \texttt{undummy} exits with \texttt{r(459)} rather than silently collapsing a respondent who chose two options into one, which is the pre-flight to run on any Qualtrics multi-select. Second, category numbering follows sort order, so pass \texttt{varnames} whenever the labels matter; it maps each category by the value its dummy carries and stays correct whether the source dummies are numeric or string.
\end{kaobox}

A note on ordering saves a later surprise: for numeric dummies the first-listed variable becomes category~1, but the order reverses for string dummies, so lean on \texttt{varnames} rather than on position; the recombined variable's labels then stay stable across platforms that export dummies as \texttt{"0"}/\texttt{"1"} strings and platforms that export them as numbers. A quick cross-check closes the loop: tabulate the recombined variable against the platform dashboard's own topline for that question; when the two disagree, the export, not the dashboard, usually dropped a category.

\section{Watching response rates while the survey is alive}\label{sec:monitor}
A survey you can only assess after it closes is a survey you cannot steer, so the highest-value move in this chapter is watching response while there is still time to act. A scheduled do-file, run daily by the operating system, pulls the current responses, counts them against the enrollment roster to get a response rate by site, and flags sites drifting below target.\marginnote{On macOS or Linux this is a \texttt{cron} entry calling \texttt{stata -b do monitor.do}; on Windows it is a Task Scheduler action. Have the do-file set a non-zero exit code when a site trips a limit, so the scheduler itself can fire the alert email.} Whether a dip is noise or news is exactly the question a control chart\index{control chart} answers, so we borrow its logic: a centerline at the expected rate and limits a few standard deviations out, with points outside the limits read as alarms rather than worries.

The chart only earns its keep if the response to it was decided before launch. In the week before the survey opens, sit with the program team and fix three things in writing: the target rate and the $k$ that sets the limits, the cadence of the pull, and the response plan: who calls whom when a site trips the lower limit, and within how many days. The cadence follows the decision cadence, not the data: if the team meets Tuesday mornings and reminders go out Wednesdays, a Monday-night pull is worth more than a daily feed nobody reads. Settle the threshold in advance and an alarm is a trigger; settle it after the field opens and it is a debate, because a limit negotiated after site~3 has already tripped it will be negotiated in site~3's favor.

\begin{vizcallout}[Visualization to build: the response-rate control chart]
Plot each site's cumulative response rate over the field period, with three reference lines via \texttt{yline()}: the target centerline and upper and lower control limits. Color out-of-limit points as alarms. A small-multiples version (\texttt{by(site)}) turns a forty-site survey into one scannable dashboard, so a program manager sees at a glance which sites need a phone call this week.
\end{vizcallout}

To make the dashboard concrete we simulate a small evaluation exactly as it would arrive from the platform: eight sites, six weekly pulls each, with a weekly response rate that centers on a target of 70 percent but wanders as reminders go out and rosters churn. The seed is fixed so the figure rebuilds identically:
{\footnotesize
\begin{verbatim}
set seed 20260704
set obs 48
egen site = seq(), block(6)      // 8 sites
bysort site: gen week = _n       // 6 weeks each
gen rate = 70 + rnormal(0, 6)    // target 70%
replace rate = rate - 14 if inlist(site,3,6) & week>=4
replace rate = min(max(rate,0),100)
\end{verbatim}
}
The two edits after the draw are the interesting part: sites~3 and~6 are pushed fourteen points down from week~4 on, so the simulated data carries two genuine laggards for the chart to catch. Everything else is noise around the target, which is what a healthy site looks like.

The control limits come from the process itself, not from a textbook constant. We set the centerline at the 70 percent target and the limits at two standard deviations of the observed weekly rates, then flag any site-week that falls below the lower limit:
{\footnotesize
\begin{verbatim}
summarize rate
local cl = 70
local sd = r(sd)
local lcl = `cl' - 2*`sd'
local ucl = `cl' + 2*`sd'
gen byte alarm = rate < `lcl'
\end{verbatim}
}
Those five locals are one small piece of statistical machinery. The centerline sits at the target, the limits stand a fixed number of standard deviations out, and a site-week trips an alarm only when it falls below the lower limit:
\[
  \text{CL}=\mu,\qquad
  \text{LCL}=\mu-k\sigma,\qquad
  \text{UCL}=\mu+k\sigma,
\]
\[
  \text{alarm}_{it}=\mathbf{1}\!\left\{\,r_{it}<\text{LCL}\,\right\}.
\]
Here $\mu$ is the target response rate (the centerline, \texttt{cl}$=70$), $\sigma$ the standard deviation of the observed weekly rates (\texttt{sd}), $k$ the number of standard deviations to the limits ($k=2$ here), $r_{it}$ the rate for site~$i$ in week~$t$, and $\mathbf{1}\{\cdot\}$ the indicator that turns on when the rate falls below the lower limit. Reading the limits in policy units is the whole point: a site-week below the lower line is not \enquote{a bit low,} it is far enough below target that noise alone rarely explains it, and that is the moment a worry becomes a phone call. With the observed weekly rates the lower limit lands at 53.7 percent, and the flagged points print as a plain alert list a program manager can act on Monday morning:
{\footnotesize
\begin{verbatim}
    site   week   rate    lcl=53.7
    ----   ----   -----   --------
       3      4   49.66   ALARM
       3      6   44.95   ALARM
       6      4   53.46   ALARM
       6      6   51.09   ALARM
    ----   ----   -----   --------
    FLAGGED_SITEWEEKS=4   FLAGGED_SITES: 3 6
\end{verbatim}
}
This makes sense: the alarm list names sites~3 and~6 and no others, which is exactly where we seeded the drop, and it names them from week~4 on, the week the drop began, so a manager watching this feed would have made two calls two weeks before the survey closed rather than reading the shortfall in the final report. That is the whole argument of the chapter in four rows: the automated pull of the previous section is what puts these two sites on a screen while there is still time to act.

An alarm names a site; it does not name a cause, and the cause decides the remedy. Three explanations cover most dips. Instrument fatigue shows up as a synchronized sag, every panel drifting down in the same field week, and its remedy is a shorter reminder or a trimmed instrument, not a phone call to one site. Site turnover shows up as one panel falling off a cliff, and a single call to the site liaison, asking whether the coordinator who championed the survey left, usually confirms it. Seasonality shows up when the dip lands on the same calendar week as last wave's, spring break or end-of-quarter, which the prior wave's monitoring output confirms in a minute. The working rule is to spend one day sorting a dip into one of these bins before escalating, because a misdiagnosed alarm spends goodwill the next one will need.

The command that draws the dashboard is a small-multiples control chart, one panel per site, with the centerline and both limits carried as reference lines:
{\footnotesize
\begin{verbatim}
twoway (line rate week, sort)                        ///
       (scatter rate week if alarm,                  ///
            mcolor(cranberry) msize(medium)),        ///
    by(site, note("") title("Weekly response rate"   ///
        " by site (simulated)", size(medium))        ///
        legend(off) rows(2))                         ///
    yline(`cl', lp(dash)) yline(`lcl', lp(dot))      ///
    yline(`ucl', lp(dot)) ytitle("Response rate (%)")///
    xtitle("Field week")
graph export "$figures/ch05_monitor.png", ///
    replace width(2400)
\end{verbatim}
}

\begin{figure*}[htbp]
\centering
\includegraphics[width=\linewidth]{ch05_monitor.png}
\caption[Small-multiples response-rate control chart]{Simulated weekly response rates for eight survey sites over a six-week field period, built by \texttt{ch05\_monitoring.do} with \texttt{set seed 20260704}. Each panel is one site; the dashed line is the 70 percent target and the dotted lines are two-standard-deviation control limits. Cranberry markers flag site-weeks below the lower limit. Sites~3 and~6 were seeded to drift down from week~4, and the chart isolates exactly those two panels, turning a forty-site survey into one scannable dashboard.}
\label{fig:ch05monitor}
\end{figure*}

Of everything in this book, the monitor is the piece a program team can act on most directly, and the team, not the analyst, is its audience. A mid-course correction, another reminder email, a call to a lagging site, happens only if someone can see mid-course, and seeing mid-course requires the automated pull of the previous section. Monitoring is where replicable data access becomes an operational tool.

One monitoring dataset feeds three products; the craft is cutting it three ways rather than sending everyone the same chart. Field staff need the names-to-call list: merge the alarm sites back against the enrollment roster, keep the nonrespondents, and hand each coordinator a short list with contact columns, because \enquote{site~3 is below the limit} is not actionable at a site and \enquote{these eleven enrollees have not responded} is. The director needs the trend and the alert, the small-multiples chart plus the four-row alarm list. The funder sees none of the mid-field churn, only the final response-rate narrative: the closing rate, the two sites that lagged, and the calls that recovered them, which presents monitoring as stewardship rather than trouble. The cut is chosen by what its reader can act on.

The last piece is the part no chart shows: something has to run the do-file every morning without anyone remembering to. On macOS or Linux one \texttt{cron} line does it, and the equivalent Windows one-liner registers a Task Scheduler job (display-only here; the paths are yours):
{\footnotesize
\begin{verbatim}
# crontab -e  (runs 7:00 every weekday, logs output)
0 7 * * 1-5 /usr/local/bin/stata-mp -b do \
    /srv/eval/ch05_monitoring.do >> /srv/eval/monitor.log 2>&1

# Windows PowerShell equivalent (Task Scheduler):
schtasks /create /tn "SurveyMonitor" /tr ^
  "stata-mp -b do C:\eval\ch05_monitoring.do" ^
  /sc weekly /d MON,TUE,WED,THU,FRI /st 07:00
\end{verbatim}
}
Without the scheduler the promise stays aspirational; with it, the do-file's non-zero exit code on an alarm lets \texttt{cron} itself mail the alert, so the whole loop, pull to phone call, runs before anyone opens a laptop.

\begin{figure*}[htbp]
\centering
\begin{tikzpicture}[node distance=7mm and 9mm]
  \node[stage] (cron) {\texttt{cron} 7:00 weekday};
  \node[stage, right=of cron] (pull) {Pull responses via API};
  \node[stage, right=of pull] (rate) {Response rate by site};
  \node[decision, right=of rate] (below) {rate\\$<$ LCL?};
  \node[io, above=8mm of below] (quiet) {No action};
  \node[io, below=12mm of below] (alert) {Alert email\\(non-zero exit)};

  \draw[flow] (cron) -- (pull);
  \draw[flow] (pull) -- (rate);
  \draw[flow] (rate) -- (below);
  \draw[flow] (below) -- node[annot,right] {no} (quiet);
  \draw[flow] (below) -- node[annot,right] {yes} (alert);

  \node[annot, below=1mm of pull] {roster + records};
  \draw[flow] (alert.west) -- ++(-1.1,0)
        |- node[annot,above,pos=0.24] {next morning} (cron.south);
\end{tikzpicture}
\caption[The scheduled monitoring loop]{The scheduled monitoring loop that runs itself. The scheduler wakes the do-file each weekday morning; the do-file pulls responses through the platform API, computes each site's response rate against the roster, and compares it to the lower control limit. A site below the limit sets a non-zero exit code that lets the scheduler mail the alert; otherwise the loop simply waits for the next morning. Nothing here waits on a human until the phone call.}
\label{fig:ch05-viz-cronloop}
\end{figure*}

\begin{kaobox}[frametitle=Ado-file Idea: \texttt{reachcheck}]
\textbf{Proposed program:} \texttt{reachcheck}. Compares the responding sample's demographics against a target-population file (Census benchmarks, or the enrollment frame) to compute representativeness ratios and flag under-covered groups while the survey is still open, so reach problems get fixed rather than footnoted.
\end{kaobox}

\section{From platform metadata to labeled variables}
A survey platform knows the text of every question and the coding of every response, and it is a waste to retype what the platform can hand you. Pulling that metadata into Stata variable characteristics with \texttt{char define}\index{char define@\texttt{char define}} lets the question wording travel with the variable, so the analyst reading \texttt{q17} six months later can recover what was actually asked without opening the instrument.\marginnote{Characteristics survive \texttt{save} and travel with the \texttt{.dta}, unlike a \texttt{note}, which is easy to overwrite. Read one back with \texttt{char list q17[]}; it is the same slot \texttt{xtset} uses, so you are storing metadata the Stata way, not bolting on your own convention.} This is accessibility built into the data itself: the meaning is attached, not stored in a separate document that will be lost.

Timing matters here: capture the metadata snapshot at launch, not at close. Platforms let an editor reword a question mid-field, and the wording your December respondents saw is recoverable only if the July pull saved it. A launch-day metadata pull, filed beside the first data pull, is the baseline for every later comparison.

Across waves, the same idea scales into an instrument history. The \texttt{surveytracker}\index{surveytracker@\texttt{surveytracker}} command logs each wave's variables, labels, and constructs into a running spreadsheet, so you can see at a glance which items are new, which changed, and which were dropped. That inventory is the backbone of the cross-wave codebook of Section~\ref{sec:cxchangelog}; without it, a multi-wave survey is not a study that extends but a pile of unrelated files.

\begin{kaobox}[frametitle=Package Integration: \texttt{surveytracker} and \texttt{validemail}]
\textbf{Integrated commands:} \texttt{surveytracker} inventories variables, labels, and constructs into a longitudinal Excel tracker to compare instruments across waves. \texttt{validemail} validates an email item in three tiers (format, live DNS/MX lookup, and disposable-provider detection), which is the small but constant chore of cleaning an open-ended contact field before it is used for follow-up.
\end{kaobox}

\section{Scales without tears}
Most survey constructs arrive as a battery of Likert items, and turning them into something reportable is repetitive enough to invite quiet errors. The \texttt{likertscale}\index{likertscale@\texttt{likertscale}} command automates the whole path: encoding text responses to numbers off the platform's value labels, reverse-coding the items that need it, computing a row-wise index, running Cronbach's alpha\index{Cronbach's alpha} \parencite{cronbach1951alpha} for reliability, and building the top-two-box\index{top-two-box} \enquote{percent agree} indicators that program boards actually read.\marginnote{Rule of thumb on alpha: below 0.70 the items are not hanging together as one scale; above 0.95 they may be redundant, several ways of asking the same question. And alpha climbs mechanically with item count, so a high value on a 30-item battery proves less than it looks.} Do this by hand across thirty items and a miscoded reverse item slips through; do it with one command and the scale rebuilds identically every wave.

To show the output a program board actually receives, we simulate a five-item coaching-satisfaction battery for the same eight sites, drawing correlated Likert responses so the items behave like one scale, then reduce each item to its top-two-box percent agree (the share answering \enquote{agree} or \enquote{strongly agree}) by site:
{\footnotesize
\begin{verbatim}
set seed 20260704
* latent site quality drives 5 correlated items,
* each collapsed to top-two-box agree (=4 or 5)
alpha q1 q2 q3 q4 q5
\end{verbatim}
}
The battery returns a scale reliability (Cronbach's alpha) of \textbf{0.91} on the pooled data, comfortably above the 0.70 floor and below the 0.95 redundancy ceiling, so the five items are measuring one coaching-satisfaction construct rather than five unrelated things. The number \texttt{alpha} reports is
\[
  \alpha=\frac{K}{K-1}\left(1-\frac{\sum_{j=1}^{K}\sigma_j^2}{\sigma_T^2}\right),
\]
where $K$ is the number of items, $\sigma_j^2$ the variance of item~$j$, and $\sigma_T^2$ the variance of the summed scale. The formula makes the margin note's warning concrete: the leading $K/(K-1)$ factor and the item variances in the sum both grow with $K$, so alpha rises with item count even when each new item adds little, which is why a high value on a thirty-item battery proves less than the same value on five. That single number licenses collapsing the battery into one row of percent-agree figures per site:

\begin{table*}[htbp]
\centering
\caption[Percent agree by site, simulated coaching battery]{Top-two-box percent agree by site for a simulated five-item coaching-satisfaction battery ($n=40$ per site, \texttt{set seed 20260704}). Pooled scale reliability is Cronbach's $\alpha=0.91$. Sites~3 and~6, the response-rate laggards of Figure~\ref{fig:ch05monitor}, are also far the lowest on satisfaction, which is the pattern a program board reads first.}
\label{tab:ch05agree}
{\footnotesize
\begin{tabular}{lcccccc}
\hline
Site & Q1 & Q2 & Q3 & Q4 & Q5 & Scale \\
     & clear & useful & respected & applied & recommend & mean \\
\hline
1 & 70.0 & 65.0 & 77.5 & 55.0 & 62.5 & 3.94 \\
2 & 57.5 & 50.0 & 60.0 & 37.5 & 52.5 & 3.54 \\
3 &  7.5 & 17.5 & 22.5 & 12.5 & 12.5 & 2.25 \\
4 & 65.0 & 60.0 & 72.5 & 37.5 & 65.0 & 3.71 \\
5 & 65.0 & 57.5 & 75.0 & 42.5 & 55.0 & 3.72 \\
6 &  5.0 &  5.0 & 10.0 &  7.5 &  7.5 & 2.02 \\
7 & 67.5 & 65.0 & 82.5 & 47.5 & 70.0 & 3.85 \\
8 & 52.5 & 52.5 & 60.0 & 37.5 & 47.5 & 3.39 \\
\hline
All & 48.8 & 46.6 & 57.5 & 34.7 & 46.6 & 3.30 \\
\hline
\end{tabular}
}
\end{table*}

Read Table~\ref{tab:ch05agree} across a row and the item ordering is stable everywhere: \enquote{felt respected} scores highest and \enquote{applied it on the job} lowest at every site, an honest signal that participants value the coaching relationship more than they act on it. Read it down the columns and sites~3 and~6 collapse to a fifth or less agreeing where the other sites clear a half to two-thirds, the same two sites the control chart flagged for slow response. That the two diagnostics point at the same sites is not a coincidence to celebrate but a lead to chase: low response and low satisfaction travel together, and the board now has both a number and a place to send help. Percent agree matters because it is the lingua franca of program audiences: \enquote{78 percent of participants agreed the coaching helped} lands where a mean of 4.1 on a five-point scale does not. Computing it consistently, from the same rule every time, is accessibility and replicability at once. The table also has more than one reader, so hold the cuts apart: the pooled bottom row and the alpha travel to the funder, the site rows go to the director, who can pair them with the control chart, and the item-level diagnostics, inter-item correlations and alpha-if-dropped, stay with the analyst.

\begin{kaobox}[frametitle=Package Integration: \texttt{likertscale}]
\textbf{Integrated command:} \texttt{likertscale}. Encodes Likert items from value-label metadata, computes row-wise indices and Cronbach's alpha, and generates collapsed top-two-box percent-agree variables with clear labels, so scale construction is one auditable step.
\end{kaobox}

\pagelayout{wide}
\section{Who did not answer, and does it matter}
A response rate is not a verdict on its own; the real question is whether the people who answered differ from the people who did not. We separate two failures: unit nonresponse\index{nonresponse!unit}, when a sampled person answers nothing, and item nonresponse\index{nonresponse!item}, when a respondent skips particular questions. The \texttt{nonresponse}\index{nonresponse@\texttt{nonresponse}} command compares respondents to the sampling frame and models the probability of responding \parencite{groves2008nonresponse}, so \enquote{can we trust a 22 percent response rate} gets a diagnostic answer instead of a shrug. The \texttt{loebias}\index{loebias@\texttt{loebias}} command adds a level-of-effort\index{level-of-effort} check, tracing whether estimates stabilize as later contact attempts bring in reluctant respondents, which is the closest thing to a window on the people you never reached. The level-of-effort logic assumes late responders resemble nonresponders, which holds better for reluctance than for unreachability, so pair the trend with a frame comparison before you lean on it: a dead phone number is not a reluctant respondent.

The reason a rate is not a verdict is that nonresponse rate and nonresponse bias are not the same quantity. \textcite{groves2006nonresponse} shows, across dozens of studies, that a survey's response rate is a weak predictor of the bias in its estimates, because bias depends on whether responding is correlated with the answer, not on how many people responded. A 40~percent response rate with responders who resemble nonresponders on the outcome can be less biased than an 80~percent rate where the missing fifth is exactly the group whose outcome differs. That is why the \texttt{nonresponse} command models who responds rather than only counting how many did, and it is why the level-of-effort trend matters: both are attempts to see the correlation, not just the count. Where the platform records it, the contact history itself is evidence. \textcite{kreuter2013paradata} makes the case for paradata\index{paradata}, the process trail of call attempts, timestamps, and reminder sends, as a lever on total survey error, and the level-of-effort check is paradata put to work: it reads the order in which responses arrived to ask whether the late, reluctant respondents are pulling the estimate.

Sequence the diagnostics the way the questions arrive. Mid-field, Section~\ref{sec:monitor} asks why the line dips; after close, this section asks whether the dip mattered, and the frame comparison comes first because it is cheap and decisive. If respondents match the frame on the characteristics that drive the outcome, a modest rate is survivable; if they do not, the gap names the group the weighting has to recover. The level-of-effort check comes second, as corroboration rather than verdict.

This matters because representativeness is the difference between \enquote{our respondents} and \enquote{our participants}, and a funder quotes the second. The diagnostics here hand off directly to the weighting of Chapter~\ref{ch:trust}, where the imbalance they find gets corrected.

\begin{kaobox}[frametitle=Package Integration: \texttt{nonresponse} and \texttt{loebias}]
\textbf{Integrated commands:} \texttt{nonresponse} compares respondents to the sampling frame with tests and logistic models and estimates raking weights \parencite{deming1940raking}; \texttt{loebias} runs cumulative level-of-effort sensitivity tests to see whether estimates stabilize as call attempts accumulate.
\end{kaobox}

\section{The cross-wave codebook}\label{sec:cxchangelog}
A survey that runs for years accumulates a quiet liability: questions get reworded, response options get added, items get dropped, and by wave nine nobody can say for certain what changed when. Reconstructing that history by hand, comparing spreadsheets wave against wave, is an afternoon's work every cycle and a source of error each time. The \texttt{cxchangelog}\index{cxchangelog@\texttt{cxchangelog}} tool regenerates the item-changes inventory automatically from a long-format crosswalk of the instrument, assigning each concept a lifecycle code (added, same, reworded, re-optioned, removed, or not asked) derived from wave-to-wave differences in wording and options.

Keeping that history is instrument stewardship. It is the measurement side of improvement science: a plan-do-study-act cycle can compare this quarter's number to last quarter's only if the instrument held still, and a question silently reworded between cycles turns an apparent improvement into an artifact \parencite{bryk2015improve}. It also keeps the utilization-focused reporting of Chapter~\ref{ch:embedded} honest across waves \parencite{patton2008utilization}: the user acting on a trend line deserves to know whether the trend reflects the program or the question.

The codebook also serves readers beyond the analyst who runs it. The per-wave summary line, one added, one removed, one reworded, is the sentence a funder report's methods section needs; the items-by-wave sheet is what the analyst opens when a trend jumps between waves, since the first suspect for a jump is the instrument, not the program. A habit that pays: add each concept's crosswalk row when the item is drafted, before launch, so the changelog is kept in real time rather than reconstructed at wave nine.

\begin{kaobox}[frametitle=Package Integration: \texttt{cxchangelog}]
\textbf{Integrated command:} \texttt{cxchangelog}\index{cxchangelog@\texttt{cxchangelog}} regenerates a cross-wave survey codebook from a long-format crosswalk of the instrument, one row per concept per wave. Install with the house idiom:
{\footnotesize
\begin{verbatim}
local gh "https://raw.githubusercontent.com/ericabooth"
net install cxchangelog, ///
    from("`gh'/cxchangelog-stata-public/main/") replace
\end{verbatim}
}
Point it at the crosswalk and name the columns:
{\footnotesize
\begin{verbatim}
cxchangelog using crosswalk_w3.xlsx, wave(wave)  ///
    concept(concept_id) wording(q_text)          ///
    options(options) study(study) summary        ///
    compare(crosswalk_prior.xlsx)                ///
    out(codebook_w3) replace
return list
\end{verbatim}
}
On the demo crosswalk this reports \enquote{concepts: 12\quad waves: 3} and a per-wave line of \enquote{1 added, 1 removed, 1 reworded,} leaving \texttt{r(n\_concepts)=12}, \texttt{r(n\_waves)=3}, and \texttt{r(n\_added)}, \texttt{r(n\_removed)}, and \texttt{r(n\_changes)} each equal to~1. It writes an Excel workbook with items-by-wave and options-by-wave sheets plus a per-wave summary of additions, changes, and removals, and \texttt{compare()} diffs a frozen prior vintage to surface what shifted between releases. Two behaviors to state plainly: a change in response-option text alone is not counted as a rewording (it shows in the options sheet instead), and an item that skips a wave and returns reworded counts as removed then added, not reworded.
\end{kaobox}

When a new analyst asks \enquote{did we change this question in 2024?}, the codebook answers in seconds, and the answer is the same no matter who runs it; that is extensibility across time and staff turnover. The synthetic multi-wave crosswalk used above, three waves and twelve concepts with one item reworded, one added, and one dropped, is generated by the package's own test battery, so the example ships with \texttt{cxchangelog} and reproduces from a clean install.

\subsection{The data-side codebook: \texttt{datadictionary}}\label{sec:datadictionary}
\texttt{cxchangelog} documents the instrument you intended; a second question remains, and it is the one the files answer. What did each wave's data actually deliver: which variables, under which value labels, with how much missingness, holding which example values? Picture the handoff where that question gets urgent. A staff-wellbeing survey has run three times, 2019, 2021, and 2023, the analyst who fielded the first two waves has left, and the third wave has just arrived as a \texttt{.dta} you did not build (180, 210, and 195 responses). Before anyone appends anything, someone has to establish what changed underneath: the workload item was dropped after 2019, a burnout item appeared in 2021, and in 2023 the supervisor-support item quietly moved from a four-point agreement scale to a five-point one, taking a new value-label set with it. Reconstructing that by opening each file and eyeballing \texttt{describe} output is the afternoon of error-prone work this chapter keeps trying to retire.

\texttt{datadictionary}\index{datadictionary@\texttt{datadictionary}} retires it with one call over the wave files. For every variable in every wave it harvests the type, display format, variable label, value-label set, nonmissing count, percent missing, distinct values, summary statistics, example values (top categories with counts for labeled items, sample values for strings), stored \texttt{notes}, and characteristics, including the \texttt{srctag} lineage the tools of Chapter~\ref{ch:panels} write. It then diffs the waves and reports every change it finds, and writes the whole codebook to a multi-sheet Excel workbook a program team can open, plus a machine-readable dataset a do-file can query:
{\footnotesize
\begin{verbatim}
datadictionary, files("staff_survey_w1.dta staff_survey_w2.dta
                 staff_survey_w3.dta")               ///
    wavenames("2019 2021 2023")                      ///
    excel(datadictionary_staffsurvey.xlsx)                 ///
    saving(codebook) replace
return list
\end{verbatim}
}
On the three demo waves this reports 25 variable-wave rows and six detected changes, and the \texttt{Changes} sheet in Figure~\ref{fig:datadictionaryxlsx} names each one: \texttt{q4} dropped in 2021, \texttt{q6} added in 2021, \texttt{q7} added in 2023, the \texttt{q1} label reworded in 2023, and the \texttt{q2} value-label set replaced (\texttt{agree4} to \texttt{agree5}) along with the display-format change that rode in with it. The \texttt{Missingness} sheet reads as a finding in its own right: the self-reported hours item was 7.8\,percent missing in 2019, spiked to 15.2 in the 2021 wave, and settled back to 9.2 in 2023, the kind of wave-specific nonresponse a trend analysis has to weigh before it interprets a dip. And the \texttt{Variables} excerpt shows \texttt{q2} carrying four distinct values under \texttt{agree4} in 2019 and five under \texttt{agree5} in 2023, with the top categories printed beside each so the recode is visible in the data, not just in the label names.

\begin{figure*}[htbp]
\centering
\includegraphics[width=\textwidth]{ch05_datadictionary_xlsx.png}
\caption[The datadictionary Excel codebook]{Three sheets of the Excel workbook \texttt{datadictionary} writes for the three-wave staff-wellbeing demo. \emph{Changes} lists every wave-to-wave difference it detected, one row each, with before and after values; \emph{Missingness} grids percent missing by wave, with \enquote{absent} marking variables a wave never carried; \emph{Variables} holds the full per-variable-per-wave codebook, excerpted here for \texttt{q2}, whose value-label set moves from \texttt{agree4} to \texttt{agree5} in 2023. The demo waves are generated by the package's test battery, so the workbook reproduces from a clean install.}
\label{fig:datadictionaryxlsx}
\end{figure*}

The reason \texttt{datadictionary} wants the wave \emph{files}, not the stitched panel, deserves one honest paragraph, because it is a fact about Stata rather than a preference of ours. When you append three waves into one file, one value-label set survives per variable; the append destroys the evidence that \texttt{q2} ever meant something different in 2019. So the codebook runs \emph{before} the stitch: document the waves as they arrived, then combine them (Chapter~\ref{ch:panels}). For a panel that is already stitched, the command's in-memory mode still earns its keep, \texttt{datadictionary, wave(wave)} grids per-wave presence and missingness, with one caveat the help file states plainly: within a stitched file, a variable that is all-missing in a wave and a variable that was never asked in that wave look identical, so the grid reports \enquote{absent or never answered} and only the wave files can say which.

\begin{kaobox}[frametitle=Package Integration: \texttt{datadictionary}]
\textbf{Integrated command:} \texttt{datadictionary}\index{datadictionary@\texttt{datadictionary}} builds an over-time codebook from wave files (\texttt{files()} or \texttt{folder()/pattern()}) or the dataset in memory, detects cross-wave changes (variables added or dropped, types, labels, value-label sets, formats), and exports a five-sheet Excel workbook (\texttt{Overview}, \texttt{Variables}, \texttt{ValueLabels}, \texttt{Changes}, \texttt{Missingness}) plus a \texttt{saving()} dataset. Returns \texttt{r(nvars)}, \texttt{r(nchanges)}, and \texttt{r(xlsx)} for scripting, so a control file can halt when \texttt{r(nchanges)} is not zero and someone has to look. Install with the house idiom:
{\footnotesize
\begin{verbatim}
local gh "https://raw.githubusercontent.com/ericabooth"
net install datadictionary, ///
    from("`gh'/datadictionary-stata-public/main/") replace
\end{verbatim}
}
It extends two community tools worth knowing: Roger Newson's \texttt{descsave} (SSC), which writes a variable-level description dataset, and Kishor Das's \texttt{codebookout} (SSC), which exports names, labels, and types to Excel. \texttt{datadictionary} adds the cross-wave diff, the missingness grid, example values, and the \texttt{notes}/\texttt{srctag} harvest. The three demo wave files are generated by the package's own test battery, so the worked example reproduces from a clean install.
\end{kaobox}

\begin{kaobox}[frametitle=Ado-file Idea: \texttt{surveypull}]
\textbf{Proposed program:} \texttt{surveypull}. One wrapper over the platform APIs of this chapter, with subcommands \texttt{responses} (download to a labeled \texttt{.dta}), \texttt{monitor} (counts and response rates against a roster, with an exit code a scheduler can read), and \texttt{codebook} (platform metadata to a data dictionary). Design principles: REDCap first because its API is simplest, one authentication story per platform, and every secret in an environment variable, never in the do-file. It fills the book's clearest tooling gap, since no existing package reaches Qualtrics or REDCap from Stata.
\end{kaobox}

\pagelayout{margin}
\paragraph{Where this leaves you.} You can now auto-download responses from the platforms evaluators use, watch response rates while a survey is live with a scheduled control chart that flags lagging sites by name, carry question meaning into the data, build scales and percent-agree indicators reproducibly with a reliability check behind them, diagnose nonresponse, and track both sides of a changing study across waves: \texttt{cxchangelog} for what the instrument asked, \texttt{datadictionary} for what the data delivered. These capabilities make an evaluation replicable, extensible, and actionable in its most operational form: they let the evaluator rebuild the data on command, carry an instrument across waves, and act on a lagging site while the survey is still open. Chapter~\ref{ch:panels} takes the data you have gathered, from APIs, scrapes, and surveys, and assembles it into panels you can defend. The thread through all of it is that each output is cut for the person who acts on it, and the thresholds that automate action were agreed before the field opened.

\section*{Exercises}\addcontentsline{toc}{section}{Exercises}
\begin{enumerate}
\item \emph{Try it on your own data.} Take a survey export you already hold, attach each question's wording to its variable with \texttt{char define}, and confirm the metadata stuck by reading one slot back with \texttt{char list}.
\item Rebuild the response-rate control chart from the chapter's simulation with \texttt{set seed 20260704}, then widen the limits from two standard deviations to three and report how many site-weeks still trip the lower limit.
\item The chapter's simulation seeds two sites to sag partway through the study. Predict, before you run anything, which sites the alarm list will name and from roughly which week; then run the monitoring do-file and check your prediction against the flagged site-weeks.
\item Break the roster on purpose by deleting one site's enrollment row, rerun the monitoring pull, and confirm the response-rate computation either errors out or reports the missing site rather than silently dividing by a wrong denominator.
\item Run \texttt{alpha} on the simulated five-item coaching battery, drop one item, and report whether the reliability rises or falls; then say in one sentence what that change tells you about whether the items hang together as one scale.
\item Draft the one-page pre-launch monitoring agreement for a survey you plan to field: the target rate, the $k$ for the control limits, a pull cadence matched to the team's meeting schedule, and the response plan naming who calls a lagging site and within how many days.
\end{enumerate}
